\documentclass[11pt]{article}
\usepackage[margin=2.2cm]{geometry}
\usepackage{amsmath,amssymb,mathtools}
\usepackage[shortlabels]{enumitem}
\setlist[enumerate,1]{label=\textbf{\arabic*.}, leftmargin=*, itemsep=10pt, topsep=6pt}
\usepackage{fancyhdr}
\usepackage{xcolor}

\newcommand{\R}{\mathbb{R}}
\newcommand{\E}{\mathbb{E}}
\DeclareMathOperator{\Var}{Var}
\DeclareMathOperator{\Cov}{Cov}
\newcommand{\ans}[1]{\par\noindent\textbf{\color{blue!50!black}Answer: #1}}

\pagestyle{fancy}
\fancyhf{}
\fancyhead[L]{Mathematical Finance \& Stochastic Processes}
\fancyhead[R]{Online Practice Quiz --- Solutions}
\fancyfoot[C]{\thepage}
\renewcommand{\headrulewidth}{0.4pt}

\begin{document}

\begin{center}
{\LARGE\bfseries Online Practice Quiz --- Solutions}\\[4pt]
{\small (10 multiple-choice questions, as administered; see \texttt{Practice Quiz questions/} for the original screenshots.)}
\end{center}
\vspace{6pt}
\hrule
\vspace{10pt}

\noindent\textbf{Notation.} $B=(B_t,t\ge0)$ denotes a Brownian motion; $\mathcal F_t=\sigma(B_s\mid0\le s\le t)$ denotes the filtration generated by $B$.

\vspace{10pt}

\begin{enumerate}

\item[\textbf{Q1.}] Determine $\displaystyle\lim_{t\to\infty}\E\big[\mathbf1_{\{B_t\le\sqrt2\}}\big]$.

\ans{$1/2$.} $B_t\sim N(0,t)$, so $\E[\mathbf1_{\{B_t\le\sqrt2\}}]=P(B_t\le\sqrt2)=\Phi(\sqrt2/\sqrt t)$. As $t\to\infty$, $\sqrt2/\sqrt t\to0$, so this converges to $\Phi(0)=1/2$. (The fixed level $\sqrt2$ becomes negligible compared to the growing spread $\sqrt t$ of $B_t$.)

\item[\textbf{Q2.}] Determine $\E\Big[\big(B_t+\int_0^tB_s\,dB_s\big)^2\Big]$.

\ans{$t+t^2/2$.} Recall $I_t:=\int_0^tB_s\,dB_s=\tfrac12(B_t^2-t)$ (Problem 5.6). Expand the square via the cross-term formula:
\[
\E\big[(B_t+I_t)^2\big]=\E[B_t^2]+2\E[B_tI_t]+\E[I_t^2].
\]
$\E[B_t^2]=t$. For the cross term, $\E[B_tI_t]=\tfrac12\E[B_t^3]-\tfrac t2\E[B_t]=0-0=0$ (odd Gaussian moment). For $\E[I_t^2]=\tfrac14\E[(B_t^2-t)^2]=\tfrac14\big(\E[B_t^4]-2t\E[B_t^2]+t^2\big)=\tfrac14(3t^2-2t^2+t^2)=\tfrac14\cdot2t^2=\tfrac{t^2}2$. Altogether $\E[(B_t+I_t)^2]=t+0+\tfrac{t^2}2=t+\tfrac{t^2}2$, confirming the stated answer.

\item[\textbf{Q3.}] Determine which type of process $e^{-t}B_{e^{2t}}$ is.

\ans{Centered Gaussian process with $\Cov(X_s,X_t)=e^{-(t-s)}$ for $0\le s<t$.} Write $X_t=e^{-t}B_{e^{2t}}$: a deterministic rescaling of Brownian motion at a deterministic (random) time, hence centered Gaussian with $\E[X_t]=e^{-t}\E[B_{e^{2t}}]=0$. For $0\le s<t$ (so $e^{2s}<e^{2t}$),
\[
\Cov(X_s,X_t)=e^{-s-t}\Cov(B_{e^{2s}},B_{e^{2t}})=e^{-s-t}\min(e^{2s},e^{2t})=e^{-s-t}e^{2s}=e^{s-t}=e^{-(t-s)}.
\]
This is the Lamperti transform exhibiting $X$ as a stationary Ornstein--Uhlenbeck process, cf.\ Problem 5.5(c)(ii).

\item[\textbf{Q4.}] Determine the value of $\E[B_s^2B_t^2]$ for $0\le s<t$.

\ans{$s(t-s)+3s^2$.} $(B_s,B_t)$ are jointly Gaussian, mean zero, with $\Var(B_s)=s$, $\Var(B_t)=t$, $\Cov(B_s,B_t)=s$. By Wick/Isserlis' theorem for a mean-zero Gaussian pair, $\E[B_s^2B_t^2]=\E[B_s^2]\E[B_t^2]+2\Cov(B_s,B_t)^2=st+2s^2$. Rewriting, $st+2s^2=(st-s^2)+3s^2=s(t-s)+3s^2$.

\item[\textbf{Q5.}] Let $X_t=B_t+t$. Determine which process is a martingale.

\ans{$e^{-2X_t}$.} For $\lambda\in\R$, $e^{\lambda X_t}=e^{\lambda B_t+\lambda t}=e^{\lambda B_t-\lambda^2t/2}\cdot e^{(\lambda+\lambda^2/2)t}$. The first factor is always a martingale, so the product is a martingale iff the deterministic factor is constant: $\lambda+\lambda^2/2=0\iff\lambda(1+\lambda/2)=0\iff\lambda=0$ or $\lambda=-2$. Hence $e^{-2X_t}$ is a martingale.

\item[\textbf{Q6.}] Determine the value of $\E\Big[\big(\int_0^1\mathbf1_{\{B_t\ge0\}}\,dB_t\big)^2\Big]$.

\ans{$1/2$.} By the Itô isometry and $\mathbf1_A^2=\mathbf1_A$,
\[
\E\left[\left(\int_0^1\mathbf1_{\{B_t\ge0\}}\,dB_t\right)^2\right]=\int_0^1\E\big[\mathbf1_{\{B_t\ge0\}}\big]\,dt=\int_0^1P(B_t\ge0)\,dt=\int_0^1\tfrac12\,dt=\tfrac12,
\]
using $P(B_t\ge0)=1/2$ for every $t>0$ by symmetry of $B_t\sim N(0,t)$.

\item[\textbf{Q7.}] Let $X_t=\int_0^te^s\,dB_s$. Determine the value of $\E\Big[\big(\int_0^tX_s\,dB_s\big)^2\Big]$.

\ans{$\tfrac14(e^{2t}-1)-\tfrac t2$.} By the Itô isometry, $\E[X_s^2]=\int_0^se^{2u}\,du=\tfrac12(e^{2s}-1)$. Applying the isometry again to the outer integral,
\[
\E\left[\left(\int_0^tX_s\,dB_s\right)^2\right]=\int_0^t\E[X_s^2]\,ds=\int_0^t\tfrac12(e^{2s}-1)\,ds=\Big[\tfrac{e^{2s}}4-\tfrac s2\Big]_0^t=\tfrac{e^{2t}}4-\tfrac t2-\tfrac14=\tfrac14(e^{2t}-1)-\tfrac t2.
\]

\item[\textbf{Q8.}] Let $X_t=B_t-tB_1$ for $0\le t\le1$ (the Brownian bridge pinned at $B_1$). Determine $\E[X_sX_t]$ for $0\le s\le t\le1$.

\ans{$s(1-t)$.} Since $\E[X_t]=0$ for all $t$, $\E[X_sX_t]=\Cov(X_s,X_t)$. For $s\le t$,
\[
\Cov(X_s,X_t)=\Cov(B_s,B_t)-t\Cov(B_s,B_1)-s\Cov(B_1,B_t)+st\Var(B_1)=s-ts-st+st=s-st=s(1-t).
\]

\item[\textbf{Q9.}] Determine the quadratic variation of $X_t=\int_0^tB_s\,dB_s$, that is, determine $\langle X,X\rangle_t$.

\ans{$\displaystyle\int_0^tB_s^2\,ds$.} The quadratic variation of an Itô integral $\int_0^tH_s\,dB_s$ is $\int_0^tH_s^2\,ds$; here $H_s=B_s$, so $\langle X,X\rangle_t=\int_0^tB_s^2\,ds$.

\item[\textbf{Q10.}] Determine which variable is a stopping time.

\ans{$2\cdot\mathbf1_{\{B_2>1\}}+4\cdot\mathbf1_{\{B_2\le1\}}$.} Write $A=\{B_2>1\}\in\mathcal F_2$ and $T=a$ on $A$, $T=b$ on $A^c$. By Problem 5.2, $T$ is a stopping time iff $a\wedge b\ge2$ (the time at which the defining event $A$ becomes known). Checking each option:
\begin{itemize}[nosep]
\item $2\cdot\mathbf1_A+\tfrac12\cdot\mathbf1_{A^c}$: $a\wedge b=1/2<2$. Not a stopping time.
\item $1\cdot\mathbf1_A+4\cdot\mathbf1_{A^c}$: $a\wedge b=1<2$. Not a stopping time.
\item $2\cdot\mathbf1_A+4\cdot\mathbf1_{A^c}$: $a\wedge b=2\ge2$. \textbf{A stopping time.}
\item $1\cdot\mathbf1_A+2\cdot\mathbf1_{A^c}$: $a\wedge b=1<2$. Not a stopping time.
\item $2\cdot\mathbf1_A+1\cdot\mathbf1_{A^c}$: $a\wedge b=1<2$. Not a stopping time.
\end{itemize}

\end{enumerate}

\end{document}
